%------------------------------------------------------------------------------%
\begin{question}
  Dadas as matrizes
  \[
    A = \begin{pmatrix}
      1 & 2 & 0 \\
      0 & 1 & 1 \\
      1 & 0 & 1
    \end{pmatrix}
    \quad
    \text{e}
    \quad
    B = \begin{pmatrix}
      1  & -2 & 2  \\
      1  & 1  & -1 \\
      -1 & 2  & 1
    \end{pmatrix}
  \]
  verifique se \(B\) é a matriz inversa de \(A\)
\end{question}

%------------------------------------------------------------------------------%
\begin{resolution}{Solução}
  Precisamos
  \begin{enumerate}
      \item Calcular o produto \(AB\)
      \item Calcular o produto \(BA\)
      \item Comparar os resultados obtidos com a matriz identidade \(I_3\)
  \end{enumerate}
\end{resolution}

%------------------------------------------------------------------------------%
\begin{resolution}{{Cálculo de \(AB\)}}
  \begin{align*}
    \onslide<+->{AB}
    \onslide<+->{& =
    \begin{pmatrix}
      1 & 2 & 0 \\
      0 & 1 & 1 \\
      1 & 0 & 1
    \end{pmatrix}
    \;
    \begin{pmatrix}
      1  & -2 & 2  \\
      1  & 1  & -1 \\
      -1 & 2  & 1
    \end{pmatrix}}
    \onslide<+->{\\ & =
    \begin{pmatrix}
      1\cdot1+2\cdot1+0\cdot(-1) & 1\cdot(-2)+2\cdot1+0\cdot2 & 1\cdot2+2\cdot(-1)+0\cdot1 \\
      0\cdot1+1\cdot1+1\cdot(-1) & 0\cdot(-2)+1\cdot1+1\cdot2 & 0\cdot2+1\cdot(-1)+1\cdot1 \\
      1\cdot1+0\cdot1+1\cdot(-1) & 1\cdot(-2)+0\cdot1+1\cdot2 & 1\cdot2+0\cdot(-1)+1\cdot1
    \end{pmatrix}}
    \onslide<+->{\\ & =
    \begin{pmatrix}
      1+2+0 & -2+2+0 & 2-2+0 \\
      0+1-1 & 0+1+2  & 0-1+1 \\
      1+0-1 & -2+0+2 & 2+0+1
    \end{pmatrix}}
  \end{align*}
\end{resolution}

%------------------------------------------------------------------------------%
\begin{resolution}{{Cálculo de \(AB\)}}
  \begin{align*}
    \onslide<+->{AB & =
    \begin{pmatrix}
      1+2+0 & -2+2+0 & 2-2+0 \\
      0+1-1 & 0+1+2  & 0-1+1 \\
      1+0-1 & -2+0+2 & 2+0+1
    \end{pmatrix}}
    \onslide<+->{\\ & =
    \begin{pmatrix}
      3 & 0 & 0 \\
      0 & 3 & 0 \\
      0 & 0 & 3
    \end{pmatrix}}
    \onslide<+->{\\ & =
    3I_3}
  \end{align*}
\end{resolution}

%------------------------------------------------------------------------------%
\begin{resolution}{{Cálculo de \(BA\)}}
  \begin{align*}
    \onslide<+->{BA}
    \onslide<+->{& =
    \begin{pmatrix}
      1  & -2 & 2  \\
      1  & 1  & -1 \\
      -1 & 2  & 1
    \end{pmatrix}
    \;
    \begin{pmatrix}
      1 & 2 & 0 \\
      0 & 1 & 1 \\
      1 & 0 & 1
    \end{pmatrix}}
    \onslide<+->{\\ & =
    \begin{pmatrix}
      1\cdot1+(-2)\cdot0+2\cdot1 & 1\cdot2+(-2)\cdot1+2\cdot0 & 1\cdot0+(-2)\cdot1+2\cdot1 \\
      1\cdot1+1\cdot0+(-1)\cdot1 & 1\cdot2+1\cdot1+(-1)\cdot0 & 1\cdot0+1\cdot1+(-1)\cdot1 \\
      (-1)\cdot1+2\cdot0+1\cdot1 & (-1)\cdot2+2\cdot1+1\cdot0 & (-1)\cdot0+2\cdot1+1\cdot1
    \end{pmatrix}}
    \onslide<+->{\\ & =
    \begin{pmatrix}
      1+0+2  & 2-2+0  & 0-2+2 \\
      1+0-1  & 2+1+0  & 0+1-1 \\
      -1+0+1 & -2+2+0 & 0+2+1
    \end{pmatrix}}
  \end{align*}
\end{resolution}


%------------------------------------------------------------------------------%
\begin{resolution}{{Cálculo de \(BA\)}}
  \begin{align*}
    \onslide<+->{BA & =
    \begin{pmatrix}
      1+0+2  & 2-2+0  & 0-2+2 \\
      1+0-1  & 2+1+0  & 0+1-1 \\
      -1+0+1 & -2+2+0 & 0+2+1
    \end{pmatrix}}
    \onslide<+->{\\ & =
    \begin{pmatrix}
      3 & 0 & 0 \\
      0 & 3 & 0 \\
      0 & 0 & 3
    \end{pmatrix}}
    \onslide<+->{\\ & =
    3I_3}
  \end{align*}
\end{resolution}

%------------------------------------------------------------------------------%
\begin{resolution}{Solução}
  Como
  \[
    AB = 3I_3 \neq I_3
  \]
  \[
    BA = 3I_3 \neq I_3
  \]
  \pause
  $B$ não é a inversa de $A$
\end{resolution}

%------------------------------------------------------------------------------%
\begin{resolution}{Solução}
  Porém
  \[
    \pause
    A^{-1} = \frac{1}{3} B
  \]
  \pause
  pois
  \[
    \pause
    AA^{-1}
    \pause
    \;=\;
    A \frac{1}{3} B
    \pause
    \;=\;
    \frac{1}{3} AB
    \pause
    \;=\;
    \frac{1}{3} \; 3 \; I_3
    \pause
    \;=\;
    I_3
  \]
  \[
    \pause
    A^{-1}A
    \pause
    \;=\;
    \frac{1}{3} BA
    \pause
    \;=\;
    \frac{1}{3} \; 3 \; I_3
    \pause
    \;=\;
    I_3
  \]
\end{resolution}

%------------------------------------------------------------------------------%
\endinput
\subsection*{Comparação com a identidade}


\subsection*{Conclusão}

Como
\[
A \cdot B = B \cdot A = 3I_3 \neq I_3,
\]
concluímos que \(B\) \textbf{não} é a matriz inversa de \(A\).

A inversa de \(A\), caso exista, seria \(\dfrac{1}{3}B\).

\end{document}